% !PRINTED 520
\weberpage{520}
\gesperrt{1. In the case of an irreducible imprimitive equation
the roots in the individual rows can be
so designated that the partial equations
$f(x, y_\alpha) = 0, f(x, y_\beta) = 0, \ldots f(x, y_\sigma) = 0$ in the field
$\Omega(\psi)$ all receive the same group.}

If one does not choose to designate the roots in the manner
indicated, the groups will at least be
isomorphic.

We further raise the question, which is in a certain measure
answered by the converse of this theorem: \gesperrt{When does a
transitive group have an intransitive normal divisor?}

If $P$ is a transitive permutation group and $Q$ an intransitive
normal divisor of $P$, let
\[
A = \alpha, \alpha_1 \ldots \alpha_{r-1}
\]
be one of the systems whose elements are interchanged by $Q$ only
among themselves, but in such a way that the elements of $A$
are transitively connected by $Q$, that is, that by permutations
from $Q$ every $\alpha$ can pass into every other $\alpha$. Now if
by an element $\pi$ from $P$ the quantities $A$ pass into
\[
B = \beta, \beta_1 \ldots \beta_{r-1}
\]
then these elements must either all coincide with the $\alpha$,
or all be different from the $\alpha$; for
by the group $\pi^{-1} Q \pi$, which by hypothesis is identical
with $Q$, the $\beta$ are interchanged only among themselves, and indeed
in the same manner as the $\alpha$, hence transitively. If therefore only a part
of the $\alpha$ occurred in $B$, then, contrary to the hypothesis,
these $\alpha$ would be interchanged by $Q$ with other quantities not
occurring in $A$.

If the $A$ do not exhaust all the elements to which
the permutations of $P$ extend, then, since $P$
is transitive, there is a permutation $\pi$ by which $A$ is carried over
into a system $B$ entirely different from it. If $\varkappa$ is a permutation
from $Q$, then by the permutation $\pi^{-1} \varkappa \pi$, which likewise
belongs to $Q$, the same permutation is produced among the indices of the $\beta$
as by $\varkappa$ among the indices of the $\alpha$. From this
it follows that by $Q$ the $\beta$ also are permuted only
among themselves. If with $A$ and $B$ the elements are not yet exhausted,
one can form in the same way a third system $C$,
which arises from $A$ by $\pi'$ and from $B$ by $\pi^{-1}\pi'$,
%CONT>>
